\documentclass[10pt,a4paper]{article}

\usepackage[utf8]{inputenc}
\usepackage[T1]{fontenc}
\usepackage[margin=1.5cm]{geometry}
\usepackage{amsmath,amssymb}
\usepackage{xcolor}
\usepackage{listings}
\usepackage{tcolorbox}
\tcbuselibrary{breakable}
\usepackage{enumitem}

\lstset{
    language=C,
    basicstyle=\ttfamily\scriptsize,
    backgroundcolor=\color{gray!10},
    keywordstyle=\color{blue},
    commentstyle=\color{gray},
    stringstyle=\color{red},
    numbers=left,
    numberstyle=\tiny,
    breaklines=true,
    frame=single,
    tabsize=4,
    showstringspaces=false,
    inputencoding=utf8,
    extendedchars=true,
    literate={é}{{\'e}}1 {è}{{\`e}}1 {à}{{\`a}}1 {ç}{{\c{c}}}1 
             {ê}{{\^e}}1 {î}{{\^i}}1 {ô}{{\^o}}1 {û}{{\^u}}1 {ù}{{\`u}}1
}

\newtcolorbox{infobox}[1]{colback=blue!5,colframe=blue!75!black,fonttitle=\bfseries,title=#1,breakable}
\newtcolorbox{exercicebox}[1]{colback=green!5,colframe=green!50!black,fonttitle=\bfseries,title=#1,breakable}
\newtcolorbox{warningbox}[1]{colback=red!5,colframe=red!75!black,fonttitle=\bfseries,title=#1,breakable}

\title{\textbf{TP Chapitre 1 : Représentation Numérique}\\
\large Exploration Pratique avec Compilation et Débogage}
\author{Dr. El Hadji Bassirou TOURÉ\\
Département de Mathématiques et Informatique\\
Université Cheikh Anta Diop de Dakar}
\date{}

\begin{document}
\maketitle

\begin{infobox}{Prérequis}
Ce TP suppose que vous avez configuré votre environnement Docker (TP-0) et maîtrisez :
\begin{itemize}
    \item Compilation avec \texttt{gcc} et options \texttt{-Wall -Wextra -g}
    \item Débogage de base avec \texttt{gdb}
    \item Utilisation des sanitizers (\texttt{-fsanitize=undefined})
\end{itemize}
\end{infobox}

\begin{infobox}{Commandes de référence}
\begin{verbatim}
# Compilation standard
gcc -Wall -Wextra -Werror -g prog.c -o prog

# Compilation avec détection d'undefined behavior
gcc -Wall -Wextra -g -fsanitize=undefined prog.c -o prog

# Débogage
gdb ./prog
(gdb) break main
(gdb) run
(gdb) print variable
(gdb) print/x variable    # Affichage hexadécimal
(gdb) print/t variable    # Affichage binaire
\end{verbatim}
\end{infobox}

%==============================================================================
\section{Exercice 1 : Affichage Binaire et Hexadécimal}
%==============================================================================

\begin{exercicebox}{Objectif}
Créer des fonctions d'affichage pour visualiser les représentations binaire et hexadécimale des entiers.
\end{exercicebox}

\textbf{1.1} Créer un fichier \texttt{affichage.c} avec le code suivant :

\begin{lstlisting}
#include <stdio.h>
#include <stdint.h>

// Affiche un entier 8 bits en binaire
void print_binary_8(uint8_t n) {
    printf("0b");
    for (int i = 7; i >= 0; i--) {
        printf("%d", (n >> i) & 1);
    }
}

int main() {
    uint8_t a = 42;
    printf("a = %3d = ", a);
    print_binary_8(a);
    printf(" = 0x%02X\n", a);
    
    // TODO: Tester avec d'autres valeurs
    
    return 0;
}
\end{lstlisting}

\textbf{Questions :}
\begin{enumerate}[label=\alph*)]
    \item Compiler et exécuter. Vérifier que 42 = 00101010 = 0x2A.
    \item Ajouter les tests pour : 0, 127, 128, 255.
    \item Écrire une fonction \texttt{print\_binary\_32(uint32\_t n)} pour les entiers 32 bits.
    \item Utiliser GDB pour afficher la valeur de \texttt{a} avec \texttt{print/t a} et \texttt{print/x a}. Comparer avec votre fonction.
\end{enumerate}

%==============================================================================
\section{Exercice 2 : Exploration Signed vs Unsigned}
%==============================================================================

\begin{exercicebox}{Objectif}
Observer les différences d'interprétation entre types signés et non signés.
\end{exercicebox}

\textbf{2.1} Créer un fichier \texttt{signed\_unsigned.c} :

\begin{lstlisting}
#include <stdio.h>
#include <stdint.h>

int main() {
    // Même motif binaire, interprétations différentes
    uint8_t u = 200;
    int8_t  s = 200;  // Attention: conversion implicite
    
    printf("=== Même valeur initiale: 200 ===\n");
    printf("uint8_t u = %u\n", u);
    printf("int8_t  s = %d\n", s);
    
    // Afficher le motif binaire de chacun
    printf("u en binaire: ");
    for (int i = 7; i >= 0; i--) printf("%d", (u >> i) & 1);
    printf("\n");
    
    printf("s en binaire: ");
    for (int i = 7; i >= 0; i--) printf("%d", ((uint8_t)s >> i) & 1);
    printf("\n");
    
    return 0;
}
\end{lstlisting}

\textbf{Questions :}
\begin{enumerate}[label=\alph*)]
    \item Compiler avec \texttt{gcc -Wall -Wextra}. Y a-t-il un warning ? Pourquoi ?
    \item Quelle valeur affiche \texttt{s} ? Expliquer avec la formule B2T.
    \item Modifier le code pour tester les valeurs : 127, 128, 255, 0.
    \item Utiliser GDB : mettre un breakpoint après les affectations, puis \texttt{print u} et \texttt{print s}. Que montre GDB ?
\end{enumerate}

%==============================================================================
\section{Exercice 3 : Pièges de Comparaison}
%==============================================================================

\begin{exercicebox}{Objectif}
Identifier et comprendre les bugs liés aux comparaisons signed/unsigned.
\end{exercicebox}

\textbf{3.1} Créer un fichier \texttt{comparaisons.c} :

\begin{lstlisting}
#include <stdio.h>
#include <stdint.h>

int main() {
    int x = -1;
    unsigned int y = 0;
    
    printf("x = %d, y = %u\n", x, y);
    printf("x < y ? %d\n", x < y);
    
    // Prédire le résultat avant d'exécuter!
    printf("-1 < 0U ? %d\n", -1 < 0U);
    printf("-1 == 0xFFFFFFFF ? %d\n", -1 == 0xFFFFFFFF);
    
    return 0;
}
\end{lstlisting}

\textbf{Questions :}
\begin{enumerate}[label=\alph*)]
    \item Prédire les résultats AVANT de compiler. Puis vérifier.
    \item Compiler avec \texttt{gcc -Wsign-compare}. Quels warnings apparaissent ?
    \item Expliquer pourquoi \texttt{-1 < 0U} est faux (quelle conversion a lieu ?).
    \item Ajouter un test : \texttt{2147483647U > -2147483648}. Prédire puis vérifier.
    \item En GDB, utiliser \texttt{print/u x} pour voir \texttt{x} interprété comme unsigned.
\end{enumerate}

%==============================================================================
\section{Exercice 4 : Détection d'Overflow avec Sanitizer}
%==============================================================================

\begin{exercicebox}{Objectif}
Utiliser le sanitizer pour détecter les undefined behaviors liés aux overflows.
\end{exercicebox}

\textbf{4.1} Créer un fichier \texttt{overflow.c} :

\begin{lstlisting}
#include <stdio.h>
#include <stdint.h>

int main() {
    // Test 1: Overflow unsigned (comportement DÉFINI)
    printf("=== Overflow Unsigned ===\n");
    uint8_t u = 255;
    printf("u = %u\n", u);
    u = u + 1;
    printf("u + 1 = %u\n", u);
    
    // Test 2: Overflow signed (UNDEFINED BEHAVIOR!)
    printf("\n=== Overflow Signed ===\n");
    int8_t s = 127;
    printf("s = %d\n", s);
    s = s + 1;  // UB!
    printf("s + 1 = %d\n", s);
    
    return 0;
}
\end{lstlisting}

\textbf{Questions :}
\begin{enumerate}[label=\alph*)]
    \item Compiler SANS sanitizer : \texttt{gcc -g overflow.c -o overflow}. Exécuter. Noter les résultats.
    \item Compiler AVEC sanitizer : \texttt{gcc -g -fsanitize=undefined overflow.c -o overflow}. Exécuter. Que se passe-t-il ?
    \item Le sanitizer détecte-t-il l'overflow unsigned ? Pourquoi ?
    \item Ajouter un test avec \texttt{int8\_t s = -128; s = s - 1;}. Le sanitizer le détecte-t-il ?
    \item Compiler avec \texttt{-O2} (optimisation) sans sanitizer. Le résultat change-t-il ?
\end{enumerate}

%==============================================================================
\section{Exercice 5 : Exploration IEEE 754}
%==============================================================================

\begin{exercicebox}{Objectif}
Visualiser la structure interne des nombres flottants.
\end{exercicebox}

\textbf{5.1} Créer un fichier \texttt{ieee754.c} :

\begin{lstlisting}
#include <stdio.h>
#include <stdint.h>

// Union pour accéder aux bits d'un float
typedef union {
    float f;
    uint32_t u;
} Float32;

void analyze_float(float value) {
    Float32 num;
    num.f = value;
    
    // Extraire les champs
    uint32_t sign = (num.u >> 31) & 1;
    uint32_t exp = (num.u >> 23) & 0xFF;
    uint32_t frac = num.u & 0x7FFFFF;
    
    printf("Valeur: %f\n", value);
    printf("Hex: 0x%08X\n", num.u);
    printf("Signe: %u\n", sign);
    printf("Exposant brut: %u (biaisé)\n", exp);
    printf("Exposant réel: %d\n", (int)exp - 127);
    printf("Fraction: 0x%06X\n", frac);
    printf("---\n");
}

int main() {
    analyze_float(1.0f);
    analyze_float(-1.0f);
    analyze_float(12.5f);
    
    // TODO: Ajouter d'autres tests
    
    return 0;
}
\end{lstlisting}

\textbf{Questions :}
\begin{enumerate}[label=\alph*)]
    \item Compiler et exécuter. Vérifier que \texttt{1.0f} = \texttt{0x3F800000}.
    \item Ajouter l'analyse de : \texttt{0.5f}, \texttt{0.1f}, \texttt{-12.5f}.
    \item Pour \texttt{0.1f}, observer la fraction. Pourquoi n'est-elle pas "propre" ?
    \item Ajouter des tests pour les valeurs spéciales :
    \begin{itemize}
        \item \texttt{1.0f / 0.0f} (infini positif)
        \item \texttt{0.0f / 0.0f} (NaN)
        \item \texttt{-0.0f} (zéro négatif)
    \end{itemize}
    \item En GDB, afficher \texttt{num.u} avec \texttt{print/t num.u} pour voir les 32 bits.
\end{enumerate}

%==============================================================================
\section{Exercice 6 : Comparaison de Flottants}
%==============================================================================

\begin{exercicebox}{Objectif}
Comprendre pourquoi la comparaison directe de flottants est dangereuse.
\end{exercicebox}

\textbf{6.1} Créer un fichier \texttt{float\_compare.c} :

\begin{lstlisting}
#include <stdio.h>
#include <math.h>

int main() {
    float a = 0.1f + 0.2f;
    float b = 0.3f;
    
    printf("a = %.20f\n", a);
    printf("b = %.20f\n", b);
    printf("a == b ? %d\n", a == b);
    printf("a - b = %.20e\n", a - b);
    
    // Comparaison correcte avec epsilon
    float epsilon = 1e-6f;
    printf("fabsf(a-b) < epsilon ? %d\n", fabsf(a - b) < epsilon);
    
    // Test: associativité
    float x = (3.14f + 1e10f) - 1e10f;
    float y = 3.14f + (1e10f - 1e10f);
    printf("\nx = (3.14 + 1e10) - 1e10 = %f\n", x);
    printf("y = 3.14 + (1e10 - 1e10) = %f\n", y);
    printf("x == y ? %d\n", x == y);
    
    return 0;
}
\end{lstlisting}

\textbf{Questions :}
\begin{enumerate}[label=\alph*)]
    \item Compiler avec \texttt{gcc -Wall -g -lm float\_compare.c -o float\_compare}. Exécuter.
    \item Expliquer pourquoi \texttt{a != b} alors que mathématiquement $0.1 + 0.2 = 0.3$.
    \item Quelle est la différence entre \texttt{x} et \texttt{y} ? Pourquoi ?
    \item En GDB, mettre un breakpoint après le calcul de \texttt{a} et \texttt{b}, puis utiliser \texttt{print a} et \texttt{print b}. GDB montre-t-il la différence ?
\end{enumerate}

%==============================================================================
\section{Exercice 7 : Projet Mini-Convertisseur}
%==============================================================================

\begin{exercicebox}{Objectif}
Créer un programme interactif de conversion entre bases.
\end{exercicebox}

\textbf{7.1} Créer un programme \texttt{convertisseur.c} qui :

\begin{enumerate}
    \item Demande à l'utilisateur un nombre (décimal, hexadécimal avec préfixe 0x, ou binaire avec préfixe 0b)
    \item Affiche ce nombre dans les trois bases
    \item Affiche l'interprétation signed et unsigned (sur 8, 16, ou 32 bits selon la valeur)
\end{enumerate}

\textbf{Exemple d'exécution attendue :}
\begin{verbatim}
Entrez un nombre (dec, 0xHEX, ou 0bBIN): 0xFF
=== Résultat sur 8 bits ===
Décimal unsigned: 255
Décimal signed:   -1
Hexadécimal:      0xFF
Binaire:          11111111

Entrez un nombre (dec, 0xHEX, ou 0bBIN): -42
=== Résultat sur 8 bits ===
Décimal unsigned: 214
Décimal signed:   -42
Hexadécimal:      0xD6
Binaire:          11010110
\end{verbatim}

\textbf{Conseils :}
\begin{itemize}
    \item Utiliser \texttt{strtol()} avec différentes bases pour parser l'entrée
    \item Créer des fonctions séparées pour chaque type d'affichage
    \item Tester avec GDB pour vérifier les conversions intermédiaires
\end{itemize}

%==============================================================================
\section{Soumission}
%==============================================================================

\begin{infobox}{À rendre}
\begin{itemize}
    \item Tous les fichiers source (.c)
    \item Un fichier \texttt{README.md} avec :
    \begin{itemize}
        \item Réponses aux questions de chaque exercice
        \item Captures d'écran montrant l'utilisation de GDB et du sanitizer
        \item Explications des bugs/comportements observés
    \end{itemize}
\end{itemize}

\textbf{Format :} Archive \texttt{TP1\_NomPrenom.tar.gz}
\end{infobox}

\end{document}
